\section{Ca Nghiên Cứu: Khắc Phục Lỗi Cắt Ngắn Ngữ Cảnh (Truncation Defect)}

\subsection{Bản chất của lỗi suy giảm chất lượng dữ liệu âm thầm}
Trong kỹ thuật dữ liệu phục vụ học máy và AI, lỗi nguy hiểm nhất không phải là các lỗi cú pháp gây sập chương trình (\textit{runtime crashes}), mà là các lỗi chất lượng dữ liệu âm thầm (\textit{silent data quality defects}) -- nơi đường ống vẫn thực thi hoàn tất thành công nhưng dữ liệu bị suy hao giá trị thông tin mà không phát ra bất kỳ cảnh báo nào.

Trong phiên bản sơ khởi của hệ thống, thuật toán phân đoạn văn bản sử dụng phương pháp đếm từ truyền thống: mỗi đoạn chứa 500 từ tiếng Anh và trượt gối đầu 100 từ, sinh ra tổng cộng 262 phân đoạn cho 7 báo cáo Form 10-K. Tuy nhiên, một từ tiếng Anh trong các tài liệu tài chính phức tạp thường tương đương từ 1,3 đến 1,8 tokens sau khi qua bộ tách từ WordPiece/BPE. Khi tiến hành kiểm định độ dài thực tế bằng bộ tách từ của mô hình \texttt{all-MiniLM-L6-v2}, đường ống phát hiện một thực trạng nghiêm trọng: có tới \textbf{257 trong tổng số 262 phân đoạn (chiếm 98,1\%)} có độ dài vượt quá giới hạn tối đa 256 tokens của mô hình nhúng véc-tơ.

Mô hình học máy xử lý tình huống này bằng cách âm thầm cắt bỏ (\textit{truncate}) toàn bộ các tokens từ vị trí 257 trở đi trong quá trình tính toán véc-tơ. Hậu quả là hơn 50\% nội dung văn bản ở nửa sau của từng phân đoạn hoàn toàn biến mất khỏi không gian véc-tơ nhúng, tạo ra các vùng ``điểm mù ngữ cảnh'' khổng lồ tại Tầng Phục Vụ.

\subsection{Thiết kế thực nghiệm đo lường và phát hiện điểm mù}
Để đo lường định lượng mức độ thiệt hại của lỗi cắt ngắn ngữ cảnh, dự án xây dựng một bộ dữ liệu kiểm định chuẩn mực (\texttt{benchmarks/eval\_questions.json}) gồm 20 câu hỏi nghiệp vụ bao quát các rủi ro cốt lõi của 7 tập đoàn công nghệ. Nhóm câu hỏi được phân bổ có chủ đích thành hai nhóm đối xứng: 10 câu hỏi có văn bản neo (\textit{anchor text}) nằm ở nửa đầu phân đoạn cũ (\texttt{late\_in\_old\_chunk: false}, độ sâu dưới 250 tokens) và 10 câu hỏi có văn bản neo nằm sâu ở nửa sau phân đoạn cũ (\texttt{late\_in\_old\_chunk: true}, độ sâu vượt quá 256 tokens).

Hệ thống đo lường hiệu năng truy xuất thông qua chỉ số Tỷ lệ tìm trúng (Hit@$k$) và Thứ hạng nghịch đảo trung bình (Mean Reciprocal Rank -- MRR), đối sánh với đường cơ sở ngẫu nhiên Monte Carlo qua 500 vòng lặp.

\subsection{Giải pháp kỹ thuật và bằng chứng thực nghiệm phục hồi dữ liệu}
Đường ống tiến hành tái cấu trúc toàn diện giải thuật phân đoạn tại Tầng Vàng (\texttt{src/rag/chunking.py}): chuyển đổi từ bộ đếm từ thô sơ sang giải thuật phân tách nhận biết bộ tách từ (\textit{tokenizer-aware chunking}). Giải thuật tôn trọng ranh giới kết thúc câu, thiết lập kích thước mục tiêu 210 tokens, bước trượt 40 tokens và áp đặt một trần cứng tuyệt đối là 240 tokens. Toàn bộ tập dữ liệu được tái cấu trúc thành 925 phân đoạn đồng đều, được bảo vệ bởi hợp đồng dữ liệu tự động \texttt{test\_gold\_chunks\_fit\_embedding\_context\_window} trong Pytest.

\begin{table}[htbp]
\centering
\scriptsize
\setlength{\tabcolsep}{3pt}
\caption{Bằng chứng thực nghiệm đối đầu trước và sau khi khắc phục lỗi cắt ngắn ngữ cảnh}
\label{tab:truncation_recovery}
\begin{tabularx}{\textwidth}{X c c c c c}
\toprule
\textbf{Cấu Hình \& Phân Vùng Câu Hỏi} & \textbf{Ngẫu Nhiên} & \textbf{Dense (MiniLM)} & \textbf{BM25 Okapi} & \textbf{Hybrid (RRF)} & \textbf{Tốt Nhất} \\
\midrule
\multicolumn{6}{l}{\textit{1. Cấu hình cũ bị lỗi (Phân đoạn 500 từ, 262 chunks, 98,1\% vượt trần context)}} \\
Toàn bộ câu hỏi Hit@5 ($\sim$3.000t) & 3,6 / 20 & 14 / 20 & \textbf{18 / 20} & 17 / 20 & \textbf{18 / 20} (BM25) \\
Nhóm nửa sau (Điểm mù) Hit@5 & --- & 6 / 10 & \textbf{9 / 10} & 8 / 10 & \textbf{9 / 10} (BM25) \\
Nhóm nửa đầu (Bình thường) Hit@5 & --- & 8 / 10 & \textbf{9 / 10} & \textbf{9 / 10} & \textbf{9 / 10} (Đồng hạng) \\
Chỉ số MRR & 0,0849 & 0,5375 & \textbf{0,6975} & 0,6417 & \textbf{0,6975} (BM25) \\
\midrule
\multicolumn{6}{l}{\textit{2. Cấu hình mới đã sửa (Token-aware $\le$240t, 925 chunks, 0\% vượt trần context)}} \\
Toàn bộ câu hỏi Hit@15 ($\sim$3.000t) & 3,1 / 20 & \textbf{18 / 20} & 16 / 20 & \textbf{18 / 20} & \textbf{18 / 20} (Dense/Hyb) \\
\textbf{Nhóm nửa sau (Điểm mù cũ) Hit@15} & --- & \textbf{9 / 10} & 7 / 10 & \textbf{9 / 10} & \textbf{9 / 10} (\textbf{+50\%}) \\
Nhóm nửa đầu (Bình thường) Hit@15 & --- & 9 / 10 & 9 / 10 & 9 / 10 & 9 / 10 (Đồng hạng) \\
Chỉ số MRR & 0,0337 & 0,4010 & \textbf{0,4326} & 0,4297 & \textbf{0,4326} (BM25) \\
\bottomrule
\end{tabularx}
\end{table}

Bảng~\ref{tab:truncation_recovery} cung cấp bằng chứng thực nghiệm rõ ràng về hiệu quả của việc sửa chữa lỗi dữ liệu từ gốc:
\begin{itemize}
    \item \textbf{Phục hồi vượt bậc tại điểm mù véc-tơ (+50\%)}: Tỷ lệ tìm kiếm chính xác (Hit@15) của mô hình Dense trên nhóm câu hỏi bị rơi vào điểm mù cũ đã tăng vọt từ 6/10 lên 9/10 (tương ứng mức phục hồi 50\%), đưa hiệu năng tổng thể của mô hình Dense đạt 18/20 (90\%), ngang bằng với hiệu năng cao nhất của BM25.
    \item \textbf{Đánh đổi về kích thước ngữ cảnh của BM25}: Ngược lại, thuật toán BM25 bị giảm nhẹ từ 18/20 xuống 16/20 (-2 câu hỏi) do các phân đoạn ngắn hơn làm mất đi một số từ khóa ngữ cảnh mở rộng xung quanh.
\end{itemize}
Ca nghiên cứu khẳng định một nguyên lý cốt lõi trong kỹ thuật dữ liệu: chất lượng và sự am hiểu tường tận về các đặc tính biểu diễn của dữ liệu quyết định giới hạn trên của toàn bộ các hệ thống AI phân tích hạ nguồn.
